\section{Conclusion and Open Problems}
\label{sec:conclusion}

Quotient-matrix reductions and exact computational sweeps constrain the
landscape of Leontovich graphs sharply. Proposition~\ref{prop:n5} rules
out $n=5$ in the spherically symmetric family $T(x,y,z)$. The smallest
crossover threshold witnessed in the
families studied here is $n=7$; two universal $n=6$ polynomial
inequalities remain open.

We conclude with the two principal open directions highlighted by this work:

\begin{enumerate}
  \item \textbf{Near-path sufficiency.} Every known Leontovich graph is
    detected by the near-path family $E_n^{(d)}$. The near-path
    sufficiency conjecture asks whether this always holds. If it does,
    the bounded near-path sweeps in this paper become bounded all-tree
    results over the same tested target orders and source-size windows;
    removing those finite bounds would still require additional arguments.
  \item \textbf{Minimal simple targets.} The bipartite double cover of
    $\hat{T}(1,35,1,50)$ yields a simple strongly Leontovich graph on
    3,644 vertices, and the perturbed quotient
    (Section~\ref{sec:doublecover}) gives high-precision numerical evidence
    for a simple non-bipartite candidate on 6,806 vertices. Closing the
    gap to the 11-vertex frontier probed by the simple-graph sweep
    remains a computational challenge.
\end{enumerate}
